\section*{Chapitre \ref{ch:iv2:options-parity-and-greeks} --- Options~: parité et grecques}
\markboth{\protect\omTocMark{}{Réponses~: Options~: parité et grecques}}{\protect\omTocMark{}{Réponses~: Options~: parité et grecques}}

\begin{solution}{iq:iv2:options-parity-and-greeks:1}
Le gain maximal est de 7~\$, la perte maximale de 3~\$, et le seuil de rentabilité se situe à 103. C'est une vision modérément haussière~: on achète un bull call spread\index{call spread}.

Le coût net vaut $5 - 2 = 3$. Sous 100, les deux options arrivent à échéance sans valeur, et on perd 3. Au-dessus de 110, le spread rapporte la totalité des 10, donc le gain vaut $10 - 3 = 7$. Entre les deux, il rapporte $S - 100$, ce qui donne le seuil de rentabilité à 103.

\iqwrong{Croire que le call vendu ajoute du risque. Il plafonne le gain, et c'est lui qui finance une partie du call acheté.}
\end{solution}

\begin{solution}{iq:iv2:options-parity-and-greeks:2}
Le seuil de rentabilité se situe à 92 et à 108. À 105, on perd 3~\$, et à 120, on gagne 12~\$. On parie que le titre bouge fortement, dans un sens ou dans l'autre, et davantage que ce que le prix des options laisse entendre.

Le gain vaut $|S - 100|$, et le coût payé est de 8. Donc à 105, $5 - 8 = -3$, et à 120, $20 - 8 = 12$. La perte maximale est de 8, en 100.

\iqwrong{Croire que tout mouvement est gagnant. Le titre doit bouger de plus de 8~\$, le montant payé, et un mouvement de 5~\$ reste perdant.}
\end{solution}

\begin{solution}{iq:iv2:options-parity-and-greeks:3}
Le profit maximal est de 8~\$, en 100. La perte maximale est égale aux 2~\$ payés. Les seuils de rentabilité sont à 92 et à 108. Ce produit ne peut pas avoir un prix négatif, car son payoff n'est jamais négatif.

En 100, le call 90 rapporte 10, et les autres ne rapportent rien. Le payoff passe donc de 0 en 90 à 10 en 100, puis retombe à 0 en 110. Déduction faite du coût de 2, le profit vaut 8 au sommet, et il est nul en 92 et en 108.

Une position qui ne peut que rapporter doit avoir un prix positif. Sinon, on est payé pour la prendre~: c'est un arbitrage butterfly\index{butterfly}.

\iqwrong{Supposer que les deux calls vendus entraînent une perte illimitée. Au-dessus de 110, les deux calls achetés les compensent exactement.}
\end{solution}

\begin{solution}{iq:iv2:options-parity-and-greeks:4}
C'est un forward long en 100, qui se comporte comme la détention du titre. Le titre doit être à 102, car le coût de la position vaut $S - K$ à taux nul. Si le titre termine à 110, on gagne 8~\$.

À l'échéance, le call rapporte $\max(S - 100, 0)$ et le put vendu rapporte $-\max(100 - S, 0)$. Ensemble, ils rapportent $S - 100$. Par la parité put-call\index{parité put-call} avec $r = 0$, $C - P = S - K$, donc $S = 102$. Le profit en 110 vaut $10 - 2 = 8$.

\iqwrong{Croire que ce montage a un risque limité parce qu'il est composé d'options. Un put vendu porte toute la baisse du titre, et elle se cache ici.}
\end{solution}

\begin{solution}{iq:iv2:options-parity-and-greeks:5}
La borne inférieure vaut $100 - 90e^{-0.05} \approx \$14.39$. À 12~\$, on achète le call, on vend le titre à découvert et on prête le produit de la vente. On gagne au moins environ 2.39~\$ en valeur d'aujourd'hui.

Aujourd'hui, on paie 12 et on reçoit 100, donc on prête 88 à 5~\%. Cela devient $88e^{0.05} \approx 92.51$. À l'échéance, si $S > 90$, on exerce le call et on achète à 90 pour clore la vente à découvert, soit un gain d'environ 2.51. Si $S < 90$, on achète au cours $S$, pour un gain de $92.51 - S$, qui est plus élevé. Actualisé, 2.51 vaut environ 2.39.

\iqwrong{Utiliser la borne $S - K = 10$, pour laquelle 12 paraît correct. Pour un call européen avec des taux positifs, le prix d'exercice est payé plus tard, il compte donc pour sa valeur actualisée.}
\end{solution}

\begin{solution}{iq:iv2:options-parity-and-greeks:6}
Il est trop cher, car le gain maximal d'un put est 100, versé dans un an~: il vaut donc au plus $100e^{-0.05} \approx \$95.12$ aujourd'hui. On le vend à 97~\$ et on prête l'argent.

On prête 97 à 5~\% pour un an, ce qui devient environ 101.97. Le montant maximal que l'on peut devoir est 100, si le titre tombe à zéro. On gagne donc au moins 1.97 à l'échéance, ce qui correspond à l'écart de 1.88 d'aujourd'hui.

\iqwrong{Comparer le put au prix d'exercice, 100, et trouver 97~\$ acceptable. Le prix d'exercice est reçu à l'échéance.}
\end{solution}

\begin{solution}{iq:iv2:options-parity-and-greeks:7}
Un put au prix d'exercice plus élevé doit coûter au moins aussi cher, et ici il coûte moins. On achète le put 110 et on vend le put 100, ce qui rapporte 1~\$ d'avance. La pire issue est un profit de 1~\$. L'écart entre les deux puts ne peut dépasser 10, la différence des prix d'exercice.

À l'échéance, la position rapporte 0 au-dessus de 110, $110 - S$ entre 100 et 110, et 10 en dessous de 100. Elle ne rapporte donc jamais moins que zéro, et on a été payé 1~\$ pour la prendre.

\iqwrong{Penser que cela s'explique par le smile. Aucun smile ne peut rendre moins cher un put au prix d'exercice plus élevé~: cela viole la règle dans tous les modèles.}
\end{solution}

\begin{solution}{iq:iv2:options-parity-and-greeks:8}
Un call acheté a un delta positif, et un put acheté un delta négatif. Le gamma et le vega sont positifs pour les deux, et le theta est négatif pour les deux. Les grandeurs communes doivent être identiques, car par parité les deux ne diffèrent que par le titre et une obligation, sans gamma, vega ni theta (à taux nul).

La parité donne $C - P = S - K$ à taux nul. En dérivant, $\Delta_C - \Delta_P = 1$. Toutes les autres dérivées sont identiques pour les deux, puisque $S - K$ n'en a aucune. Donc quel que soit le gamma, le vega ou le theta d'un call, le put au même prix d'exercice les a également.

\iqwrong{Dire qu'un put a un gamma négatif parce que son delta est négatif. Le delta augmente avec le titre pour les deux, de $-1$ vers 0 pour un put.}
\end{solution}

\begin{solution}{iq:iv2:options-parity-and-greeks:9}
Le delta correspond à une position longue de 550 actions. On vend 550 actions pour le ramener à zéro.

Calls~: $10 \times 100 \times 0.6 = 600$. Puts vendus~: $-5 \times 100 \times (-0.3) = +150$. Titre~: $-200$. Et $600 + 150 - 200 = 550$.

\iqwrong{Oublier la double négation sur le put vendu. Un put vendu gagne quand le titre monte, il est donc long en delta.}
\end{solution}

\begin{solution}{iq:iv2:options-parity-and-greeks:10}
On vend 250 actions pour se couvrir. Après une hausse de 1~\$, le delta vaut 0.52, donc les calls représentent une position longue de 260 actions, et on vend 10 de plus.

$5 \times 100 \times 0.5 = 250$. Le nouveau delta vaut $0.5 + 0.02$. Donc $5 \times 100 \times 0.52 = 260$. La couverture vend davantage quand le titre monte, et en achète davantage quand il baisse. Être long en gamma signifie vendre haut et acheter bas.

\iqwrong{Se couvrir une fois et croire que c'est fini. Le delta évolue avec le cours, et le gamma mesure la vitesse de cette évolution.}
\end{solution}

\begin{solution}{iq:iv2:options-parity-and-greeks:11}
Avec un mouvement de 2~\$, on gagne $\tfrac12 \times 1000 \times 4 - 1500 = +\$500$. Avec un mouvement de 1~\$, on perd 1\,000~\$. Le seuil de rentabilité se situe pour un mouvement d'environ 1.73~\$.

Le gain du gamma vaut $\tfrac12\Gamma m^2$. Pour le seuil de rentabilité, $500m^2 = 1500$, donc $m = \sqrt 3$.

\iqwrong{Oublier le facteur un demi. Il est là parce que le delta gagné passe de zéro à sa pleine taille pendant le mouvement.}
\end{solution}

\begin{solution}{iq:iv2:options-parity-and-greeks:12}
On gagne 1\,100~\$. Un mouvement unique de 3 points sur la totalité du vega donnerait 1\,500~\$.

$300 \times 3 + 200 \times 1 = 1\,100$. Les surfaces ne bougent pas en parallèle, car la volatilité à court terme bouge davantage~: il faut donc un vega par tranche d'échéance.

\iqwrong{Donner un seul chiffre de vega total pour un livre. Un livre peut avoir un vega nul et perdre lourdement si les échéances proches et lointaines évoluent différemment.}
\end{solution}

\begin{solution}{iq:iv2:options-parity-and-greeks:13}
Le call gagne environ 0.47~\$, et le put perd environ 0.53~\$. Un taux plus élevé réduit la valeur actualisée du prix d'exercice. Le détenteur du call, qui paie ce prix d'exercice, gagne donc, tandis que le détenteur du put, qui le reçoit, perd.

Le rho\index{rho} d'un call vaut $KTe^{-rT}\Phi(d_2)$ par 1.00 de taux, soit ici $100 \times \Phi(-0.1) \approx 46$, ou 0.46 par point à taux nul. Pour un put, il vaut $-100 \times \Phi(0.1) \approx -54$. Sur un point complet, le rho du call croît un peu, ce qui donne 0.47 et $-0.53$. La parité donne le même total~: un point sur la valeur actualisée de 100 vaut environ 1.

\iqwrong{Penser que les taux font bouger les actions et les options de la même façon. Ici le titre est maintenu fixe. En pratique, on ne peut pas tout maintenir fixe.}
\end{solution}

\begin{solution}{iq:iv2:options-parity-and-greeks:14}
Un call américain vaut au moins la borne européenne $S - Ke^{-rT}$, qui dépasse $S - K$, c'est-à-dire ce que procure l'exercice. Le revendre est donc toujours préférable à l'exercer. Pour un put, l'exercice donne le prix d'exercice tout de suite, et la même borne joue dans le sens inverse.

Pour un call européen, la borne inférieure vaut $S - Ke^{-rT}$. Avec un $r$ positif, elle dépasse $S - K$ de l'intérêt sur le prix d'exercice. Un call américain vaut au moins autant qu'un call européen. Donc l'exercice rapporte moins que le prix de l'option.

Pour un put, la borne vaut $Ke^{-rT} - S$, inférieure au $K - S$ de l'exercice. Ainsi, si le put est assez profondément dans la monnaie, exercer plus tôt permet de toucher le prix d'exercice et ses intérêts. Par ailleurs, la valeur d'assurance d'un call disparaît pour rien si on l'exerce.

\iqwrong{«~Il faut exercer si l'option est profondément dans la monnaie.~» Il faut la revendre, ou la couvrir, et ne pas l'exercer, tant qu'il n'y a pas de dividende.}
\end{solution}

\begin{solution}{iq:iv2:options-parity-and-greeks:15}
L'écart devrait être $100 - 95.12 = 4.88$, et il vaut 3. Le call est donc bon marché par rapport au put. On achète le call, on vend le put, on vend le titre à découvert et on prête le reste. Cela rapporte environ 1.88~\$ aujourd'hui.

Acheter le call et vendre le put coûte 3. Vendre le titre à découvert rapporte 100. On prête 95.12 sur les 97, et on garde 1.88. À l'échéance, l'obligation verse 100. Le call et le put permettent ensemble d'acheter le titre à 100, quelle que soit l'issue, et on s'en sert pour clôturer la vente à découvert. Tout se compense donc à zéro, quel que soit le cours du titre.

Ce montage s'appelle un reversal\index{reversal}, une conversion effectuée dans l'autre sens, et il ne laisse subsister aucun risque sur le titre.

\iqwrong{Ne trader que les options. Sans le titre ni l'obligation, on détient une opinion, et non un arbitrage.}
\end{solution}

\begin{solution}{iq:iv2:options-parity-and-greeks:16}
Les frais effacent l'écart. Un écart de 1.88 sur un titre à 100~\$ représente environ 1.88~\% sur un an, si bien qu'un frais de 2~\% suffit à le couvrir et au-delà. L'écart implique donc des frais d'emprunt d'environ 1.9~\%. Il renseigne donc sur le coût de la vente à découvert, et ce n'est pas un profit facile.

Réaliser l'opération suppose de vendre le titre à découvert. Payer 2~\% sur 100 pendant un an coûte environ 2, soit plus que les 1.88 que l'on a sécurisés.

Un titre difficile à emprunter a un forward inférieur au carry habituel, et les options le révèlent par la parité. Les marchés d'options fournissent souvent la cotation la plus juste des frais d'emprunt d'un titre.

\iqwrong{Dire au recruteur qu'on a trouvé une anomalie de prix. Il faut d'abord se demander~: y a-t-il un dividende, quels sont les frais d'emprunt, et la vente à découvert peut-elle être rappelée~?}
\end{solution}

\begin{solution}{iq:iv2:options-parity-and-greeks:17}
Les bornes sont $S - K = 0$ et $S - Ke^{-rT} \approx 4.88$. Donc $0 \leq C - P \leq 4.88$.

Pour la borne supérieure, $C \leq S - Ke^{-rT} + P$. Comme un call américain sur un titre sans dividende vaut autant qu'un call européen, la parité européenne donne $C = S - Ke^{-rT} + P_E$, et $P_E \leq P$. Donc $C - P \leq S - Ke^{-rT}$.

Pour la borne inférieure, il faut acheter le call, détenir $K$ en liquidités, vendre le put et vendre le titre à découvert. Si le put est exercé contre la position à tout moment, on paie $K$ avec les liquidités, on reçoit une action et on clôture la vente à découvert, et on garde encore le call et les intérêts. Si le put n'est jamais exercé, alors à l'échéance le call et le put vendu versent ensemble $S_T - K$, et, compte tenu de la vente à découvert du titre et de la trésorerie, il ne reste que les intérêts. La valeur finale n'est donc jamais négative, et la position ne peut pas coûter moins que zéro~: $C + K - P - S \geq 0$.

\iqwrong{Utiliser l'égalité. Un put américain vendu peut être exercé contre la position à tout moment~; on ne peut donc pas rendre exacte l'opération de parité, et on obtient un encadrement.}
\end{solution}

\begin{solution}{iq:iv2:options-parity-and-greeks:18}
Il verse exactement 20~\$, quelle que soit la clôture du titre. C'est donc une obligation zéro-coupon d'un an. À 18.70~\$, il implique un taux d'environ 6.72~\%, si bien qu'on emprunte à 5~\% pour l'acheter.

Les deux premières jambes forment un forward long à 90, et les deux dernières un forward court à 110. Ensemble, elles versent $(S - 90) - (S - 110) = 20$. Donc $18.70e^{r} = 20$ donne $r = \ln(20/18.70) \approx 6.72~\%$. À 5~\%, la fair value du box vaut $20e^{-0.05} \approx 19.02$.

C'est un box spread\index{box spread}~: un moyen de prêter ou d'emprunter par le biais d'options. C'est donc un moyen de mesurer le taux implicite dans les options.

\iqwrong{Le traiter comme gratuit. Cela suppose de pouvoir tenir la position jusqu'à l'échéance, et les options américaines, les appels de marge et le crédit de la chambre de compensation entrent tous en jeu.}
\end{solution}

\begin{solution}{iq:iv2:options-parity-and-greeks:19}
Oui, sur un titre sans dividende, il s'agit d'un arbitrage. Le call le plus long doit valoir au moins autant que le plus court. Avec un gros dividende entre les deux, ce n'est plus le cas~: les détenteurs du call le plus long passent à côté du dividende. Pour réaliser l'opération, on vend le plus court à 4~\$, on achète le plus long à 3.50~\$, et on encaisse 0.50~\$.

À la première échéance, si le call vendu est dans la monnaie, on doit $S - 100$. Le call acheté a encore deux mois, et il vaut au moins $S - 100e^{-r\tau}$, ce qui dépasse $S - 100$. On peut donc le vendre, et on est toujours couvert.

Avec un dividende entre les deux échéances, la baisse de cours fait baisser le call le plus long, dont la valeur peut être plus faible.

\iqwrong{Dire qu'une option plus longue coûte toujours davantage. Cela suppose qu'il n'y ait aucun dividende entre les deux, et, pour un put européen, la règle peut aussi échouer lorsque les taux sont élevés.}
\end{solution}

\begin{solution}{iq:iv2:options-parity-and-greeks:20}
Le gamma vaut environ 0.0198 avec un an restant et 0.317 avec un jour restant, soit environ 16 fois plus. Avec un jour restant, le gamma est divisé par deux à environ 1.48~\$ du prix d'exercice. Les options vendues à courte échéance portent donc un risque important, qui évolue très vite.

$\Gamma = \varphi(d_1)/(S\sigma\sqrt T)$. Avec un an~: $\varphi(0.1)/20$. Avec un jour~: $\sigma\sqrt T = 0.2/\sqrt{252} \approx 0.0126$, donc $\Gamma \approx 0.399/1.26 \approx 0.317$. Le gamma est $\varphi$ de la distance standardisée, il est donc divisé par deux lorsque $d = \sqrt{2\ln 2} \approx 1.18$, c'est-à-dire à environ 1.18 écarts-types journaliers, soit 1.48~\$.

Une position vendeuse en gamma près d'un prix d'exercice en fin de vie représente donc une exposition importante. Sur une plage de 3~\$ centrée sur le prix d'exercice, le delta passe d'environ 0.12 à environ 0.88.

\iqwrong{Penser que le risque baisse à mesure que l'option approche de l'échéance parce qu'elle devient bon marché. Sa valeur baisse, mais son gamma, près du prix d'exercice, augmente.}
\end{solution}

\begin{solution}{iq:iv2:options-parity-and-greeks:21}
Elle vaut environ 45.13~\$. Elle gagne environ 1.3 centimes avec un jour de plus, soit environ 4.74~\$ par an. Le theta peut être positif de cette façon lorsqu'un put européen est profondément dans la monnaie avec un taux positif, ou lorsqu'un call profondément dans la monnaie présente un fort rendement de dividende.

Le put verse presque sûrement 100 moins le titre, et il se comporte donc presque comme une obligation~: $100e^{-0.05} - 50 \approx 45.12$. Au fil du temps, le prix d'exercice actualisé croît, de $rKe^{-rT} \approx 4.76$ par an, et c'est à peu près le theta.

\iqwrong{«~Le theta est toujours négatif.~» C'est vrai pour une option à taux nul. Un put européen profondément dans la monnaie est surtout une obligation, et il rapporte des intérêts comme une obligation.}
\end{solution}

\begin{solution}{iq:iv2:options-parity-and-greeks:22}
Les vegas valent environ 0.19, 0.40 et 0.29 par point, pour des prix d'exercice de 80, 100 et 120. Ils ne sont pas symétriques autour de 100, pour deux raisons~: le vega culmine pour $d_1 = 0$, au-dessus du cours, et un écart de 20 points est plus grand en termes logarithmiques en dessous qu'au-dessus.

Le vega vaut $S\varphi(d_1)\sqrt T$. Les valeurs de $d_1$ sont environ 1.22, 0.10 et $-0.81$. Le sommet, $d_1 = 0$, se situe au prix d'exercice $Se^{\sigma^2T/2} \approx 102$. Et le prix d'exercice de 120 est plus proche en logarithmes~: $\ln(120/100) = 0.18$, contre $\ln(80/100) = -0.22$.

Un prix d'exercice supérieur de 20~\% a donc plus de vega qu'un prix inférieur de 20~\%. Cela compte lorsqu'on monte une opération neutre en vega.

\iqwrong{Supposer que des écarts égaux en dollars ont des grecques égales. La loi normale s'exprime en termes logarithmiques.}
\end{solution}

\begin{solution}{iq:iv2:options-parity-and-greeks:23}
Le delta vaut $\Phi(0.2) \approx 0.579$, et la probabilité vaut $\Phi(-0.2) \approx 0.421$. Le trader se trompe d'environ 0.16. L'affirmation n'est juste que lorsque la volatilité totale est faible.

Ici, $d_1 = \sigma\sqrt T/2 = 0.2$ et $d_2 = -0.2$. Leur écart vaut $\sigma\sqrt T$. Ainsi, la probabilité que l'option finisse dans la monnaie est $\Phi(d_2)$, sous la probabilité risque-neutre, et le delta vaut $\Phi(d_1)$.

Le delta est la probabilité sous la mesure qui prend le titre comme numéraire. Cette mesure accorde plus de poids aux prix élevés, ce qui explique qu'il soit plus élevé.

\iqwrong{Dire que le delta est la probabilité de finir dans la monnaie. Les deux sont proches pour les options à courte échéance et à faible volatilité, et c'est de là que vient l'habitude.}
\end{solution}

\begin{solution}{iq:iv2:options-parity-and-greeks:24}
Le delta du put vaut environ $-0.43$. Sans rendement de dividende, il vaudrait $-0.45$.

Avec un rendement de dividende $q$, la parité s'écrit $C - P = Se^{-qT} - Ke^{-rT}$, et en dérivant, on obtient $\Delta_C - \Delta_P = e^{-qT}$. Donc $\Delta_P = 0.55 - e^{-0.02} \approx 0.55 - 0.9802 = -0.4302$.

\iqwrong{Utiliser $\Delta_P = \Delta_C - 1$ dans tous les cas. C'est le cas sans dividende, et l'erreur atteint 2~\% ici.}
\end{solution}

\begin{solution}{iq:iv2:options-parity-and-greeks:25}
On vend 0.12~action à 103, et on les rachète à 100, soit un gain de 0.36~\$ par straddle. La formule donne $2 \times \tfrac12 \times 0.04 \times 9 = 0.36$. Un theta journalier de $-0.36$ ramènerait au point mort.

À 100, le delta est nul. À 103, le delta vaut $0.04 \times 3 = 0.12$, donc on vend 0.12 à 103. De retour à 100, le delta est de nouveau nul, donc on rachète 0.12 à 100. Cela fait $0.12 \times 3 = 0.36$. Le gain vient donc de la vente au plus haut et de l'achat au plus bas. C'est le gamma scalping\index{gamma scalping}.

\iqwrong{Penser qu'un titre inchangé signifie aucun gain dans la journée. Le gain du hedger dépend de l'amplitude des mouvements du titre entre-temps.}
\end{solution}

\begin{solution}{iq:iv2:options-parity-and-greeks:26}
On aurait dû gagner 50~\$~: $+400$ au titre du delta, $+100$ au titre du gamma, $-300$ au titre du vega et $-150$ au titre du theta. On a gagné 120~\$, donc 70~\$ restent inexpliqués.

Delta~: $200 \times 2$. Gamma~: $\tfrac12 \times 50 \times 4$. Vega~: $300 \times (-1)$. Theta~: $-150$.

Les 70~\$ peuvent provenir d'une variation du skew, d'un mouvement du vega avec le sous-jacent (vanna), d'un chemin intrajournalier si l'on a rééquilibré la couverture, de dividendes, d'une erreur de données, ou de transactions réalisées dans la journée à un bon prix. Avant de faire confiance au modèle, il faut déterminer laquelle de ces causes est en jeu. Un montant inexpliqué qui persiste est un signal d'alerte.

\iqwrong{Ignorer un petit montant inexpliqué au motif que le total est proche. Les erreurs peuvent se compenser, et un gain inexpliqué constitue un signal d'alerte au même titre qu'une perte.}
\end{solution}

\begin{solution}{iq:iv2:options-parity-and-greeks:27}
À l'échéance, on perd 9.50~\$ à 80, on gagne 0.50~\$ à 100 et on gagne 10.50~\$ à 120. Le delta vaut environ 0.62 et le vega environ 0.06 par point. On gagne environ 0.68~\$ pour une hausse d'un point de la volatilité implicite à 110 et une baisse d'un point de celle à 90.

Le crédit net est de 0.50. À 80, le put coûte 10. À 120, le call rapporte 10. Le delta vaut $0.35 + 0.27$, et le vega $0.37 - 0.31$. Une hausse de la volatilité du call rapporte 0.37, et une baisse de celle du put vendu rapporte 0.31. On détient donc un risk reversal\index{risk reversal}, qui est surtout un pari sur la pente du smile, et sur la direction.

\iqwrong{Qualifier cette position de gratuite parce qu'elle ne coûte rien ou qu'elle rapporte. Un crédit ne signifie pas l'absence de risque, puisque l'on supporte toute la baisse en dessous de 90.}
\end{solution}

\begin{solution}{iq:iv2:options-parity-and-greeks:28}
Le gain maximal et la perte maximale valent tous deux 10~\$. La position équivaut à l'achat d'un call spread 90-110 et à la détention de 90~\$ en liquidités. Ce call spread doit donc coûter 10~\$.

À l'échéance, la position vaut $\min(\max(S, 90), 110)$. Cela égale $90 + \max(S - 90, 0) - \max(S - 110, 0)$. Par la parité avec des taux nuls, $S + P_{90} = C_{90} + 90$. Donc $S + P_{90} - C_{110} = 90 + C_{90} - C_{110}$. Le coût à gauche vaut $100 + 3 - 3 = 100$, donc le spread coûte 10.

\iqwrong{Qualifier un collar\index{collar} d'«~assurance gratuite~». Le coût de la protection a été payé en renonçant au gain au-dessus de 110.}
\end{solution}

\begin{solution}{iq:iv2:options-parity-and-greeks:29}
Le profit vaut $-\$12$ à 85, $+\$8$ à 105 et $+\$8$ à 120. Il équivaut à la vente d'un put 105, qui doit se négocier à 8~\$.

$S + (-\text{call}) = \min(S, 105)$ à l'échéance, plus la prime, moins les 100 payés. Par la parité avec des taux nuls, $P = C - S + K = 3 - 100 + 105 = 8$. Un put vendu à 8 rapporte $8 - 20 = -12$ à 85, et 8 au-dessus de 105.

La vente d'un call couvert revient donc à vendre un put. On conserve presque toute la baisse, et l'on renonce à la hausse, en échange d'une prime modeste.

\iqwrong{«~Revenu~». La prime est la rémunération de la prise de risque. La position est un put vendu, et en cas de forte baisse elle ressemble beaucoup à la détention de l'action.}
\end{solution}

\begin{solution}{iq:iv2:options-parity-and-greeks:30}
La différence entre le call et le put vaut environ 4.88~\$~: $(100 - 95)e^{-0.025}$.

Être long d'un call et short d'un put au même strike revient à s'engager à acheter à terme à 95. La valeur du contrat à terme à l'échéance est $F$, donc le gain est $F_T - 95$. Un contrat à terme ne coûte rien à l'entrée, donc sa juste valeur aujourd'hui est la valeur actuelle de $F - K$. D'où $C - P = (F - K)e^{-rT}$.

La version sur action fait apparaître $S - Ke^{-rT}$~: l'action doit être payée aujourd'hui, alors qu'un contrat à terme ne l'est pas. Si l'on pose $F = Se^{rT}$, les deux expressions coïncident.

\iqwrong{Utiliser $S - Ke^{-rT}$ avec le prix du contrat à terme à la place de $S$. On obtient alors $100 - 95e^{-0.025} \approx 7.35$, ce qui compte un coût de financement que l'on ne paie pas.}
\end{solution}

\begin{solution}{iq:iv2:options-parity-and-greeks:31}
Elle vaut environ 0.46~\$, soit $\Phi(d_2)$ avec $d_2 = -0.1$. C'est la limite d'un call spread serré~: $-\partial C/\partial K$. Près de l'échéance et au voisinage du strike, son delta devient énorme, et la couverture fait des sauts.

Un spread qui paie 1 au-dessus de $K + \varepsilon$, constitué de $1/(2\varepsilon)$ calls de strikes $K - \varepsilon$ et $K + \varepsilon$, tend vers le digital. Son prix vaut $-\partial_K C$. Dans Black-Scholes, cela donne $\Phi(d_2)$. Une vérification numérique avec $\varepsilon = 0.5$ donne le même résultat à trois décimales près.

Le delta vaut $\varphi(d_2)/(S\sigma\sqrt T)$, et il croît sans limite lorsque $T \to 0$, le sous-jacent étant au strike.

\iqwrong{Prendre le delta d'un call, $\Phi(d_1)$, pour le prix. On obtiendrait 0.54 ici. Un digital payant en numéraire se valorise par la probabilité risque-neutre.}
\end{solution}

\begin{solution}{iq:iv2:options-parity-and-greeks:32}
Son gamma est positif près de 100 et en dessous, négatif près de 110 et au-dessus, et il s'annule entre les deux, vers le milieu de l'intervalle. Cela tient au fait que l'on est long de gamma au strike inférieur et court au strike supérieur, et que le gamma de chaque option culmine près de son propre strike.

Dans Black-Scholes, les deux gammas sont égaux là où les deux $d_1$ sont égaux et opposés, soit en $S = \sqrt{K_1K_2}\,e^{-\sigma^2T/2}$. Avec un mois restant et une volatilité de 20~\%, cela donne environ 104.7.

Ainsi, à mesure que le sous-jacent monte dans l'intervalle, votre delta augmente d'abord, puis diminue. En rééquilibrant, on vend lors d'une hausse vers 100, et l'on achète lors d'une hausse vers 110. Le spread perd du theta près du strike inférieur, et il en encaisse près du strike supérieur.

\iqwrong{Considérer les grecques d'un spread comme fixes. Leur signe dépend du niveau du sous-jacent, si bien que la couverture peut s'inverser au cours d'une journée.}
\end{solution}

\begin{solution}{iq:iv2:options-parity-and-greeks:33}
Le call à 25~\% vaut environ 9.95~\$, et le put à 27~\% environ 10.74~\$. Ils s'écartent donc d'environ 0.79~\$ de la parité, ce qui revient à dire que les deux volatilités implicites sont distantes de 2 points. Le taux d'emprunt implicite est d'environ 0.8~\% par an.

Avec des taux nuls, la parité donne $C = P$ au strike. Un vega de 0.40 par point, sur un écart de 2 points, représente environ 0.80~\$. Il suffit d'abaisser le forward de 0.79~\$ pour que les deux options s'ajustent à une seule volatilité. Ce que l'on prend pour un écart de volatilité est donc en réalité un mauvais forward.

La plupart des systèmes résolvent donc un seul forward implicite à partir du call et du put proches de la monnaie. Ils utilisent ensuite une volatilité par strike, en général celle de l'option hors de la monnaie.

\iqwrong{Chercher à arbitrer les deux volatilités sans s'interroger sur le forward. L'arbitrage requiert l'action~: si la vente à découvert en est coûteuse, l'écart peut subsister.}
\end{solution}

\begin{solution}{iq:iv2:options-parity-and-greeks:34}
Oui, il faut exercer maintenant. L'exercice rapporte 99 aujourd'hui, alors qu'un put européen ne peut jamais valoir plus que la valeur actuelle du strike, $100e^{-0.05} \approx 95.12$. À taux nul, on n'exercerait pas, car le strike ne rapporte aucun intérêt, et l'on renoncerait à la valeur résiduelle de l'option.

Conserver le put jusqu'à l'échéance rapporte au plus 100 dans un an, ce qui vaut 95.12 aujourd'hui. Prendre 99 maintenant et le placer donne environ 104.08 dans un an. L'exercice anticipé est donc préférable au meilleur cas de la détention.

\iqwrong{Appliquer au put le résultat du call. Pour un call, attendre est un avantage~: on conserve le strike. Pour un put, attendre coûte, car on attend d'être payé.}
\end{solution}

\begin{solution}{iq:iv2:options-parity-and-greeks:35}
Avec un an restant, on détient environ 19\,850 actions. Avec un jour restant, on en détient environ 316\,600, pour un paiement de 1~million de dollars. Un dealer le valorise comme un call spread qui paie toujours au moins le digital. Cela plafonne le delta et permet de couvrir avec des options ordinaires.

La couverture vaut $\$1\,000\,000 \times \varphi(d_2)/(S\sigma\sqrt T)$. Avec un an~: $1\,000\,000 \times 0.397/20 \approx 19\,850$. Avec un jour~: environ 316\,600. Le jour de l'échéance, un mouvement de 1~\$ du titre entraîne donc une variation considérable, et il faut échanger un nombre énorme d'actions juste autour du strike.

Pour réduire le risque, on détient contre le digital un call spread de largeur $w$~: on est long de $1\,000\,000/w$ calls de strike $K - w$ et court d'autant de calls de strike $K$. Il paie l'intégralité du montant au-dessus de $K$, ce qui couvre ce que l'on doit, et un peu plus juste en dessous. Avec $w = 2$, son delta ne dépasse jamais 500\,000 actions, quelle que soit la proximité de l'échéance, alors que le delta du digital n'a pas de borne. La prime supplémentaire est le prix de ce confort.

\iqwrong{Couvrir un digital comme on couvre une option vanille, et faire confiance au delta du modèle. Le delta du modèle est énorme au voisinage du strike, et sur un marché réel on ne peut pas le négocier.}
\end{solution}

\begin{solution}{iq:iv2:options-parity-and-greeks:36}
Mesurée à la clôture, si l'on est assigné sur la totalité, on ne perd rien. Si l'on n'est pas assigné du tout, on perd 200\,000~\$. C'est le risque de pin\index{risque de pin}. Il fallait clôturer la position, ou au moins la réduire, avant l'échéance.

À la clôture, chaque action, diminuée de son call vendu, valait $100.02 - 0.02 = 100$.

Assigné~: on livre les 100\,000 actions à 100~\$, valeur qui est précisément la même, et l'on ne supporte plus aucun risque sur le titre.

Non assigné~: les calls expirent, et l'on détient toujours 100\,000 actions, désormais à 98~\$. Cela fait $-\$2$ par action, soit 200\,000~\$.

Et il faut s'attendre au mauvais scénario. Un porteur qui voit 98~\$ avant l'échéance n'exercera pas, et celui qui voit 102~\$ exercera, si bien que le choix joue contre le vendeur de l'option. Il faut clôturer la couverture et les options en fin de journée, même à un coût modeste, afin de ne pas conserver une position inconnue pendant la nuit.

\iqwrong{Qualifier la position de neutre en delta. Le dernier jour, le delta vaut soit 0, soit 100\,000 actions, et l'on découvre lequel après que le marché a bougé.}
\end{solution}

\begin{solution}{iq:iv2:options-parity-and-greeks:37}
Oui, il faut exercer. On gagne environ 0.36~\$ par rapport à l'attente. La règle consiste à exercer lorsque le dividende dépasse $K(1 - e^{-r\tau})$ augmenté de la valeur du put. Ici, $2 > 1.24 + 0.40$.

Exercer maintenant donne $110 - 100 = 10$. Si l'on attend, l'action baisse à 108. Par la parité, sans dividende supplémentaire, le call vaut $P + (S - D) - Ke^{-r\tau} = 0.40 + 108 - 98.76 = 9.64$. L'écart est donc de 0.36.

En général, on compare $S - K$ avec $P + S - D - Ke^{-r\tau}$. On exerce donc si $D > K(1 - e^{-r\tau}) + P$. On reçoit le dividende, mais l'on renonce à deux choses~: les intérêts sur le strike et l'assurance qu'apporte une option encore vivante.

L'exercice se décide donc juste avant la date de détachement. Un teneur de marché doit surveiller chaque option américaine avant chaque dividende.

\iqwrong{Exercer parce que le call est profondément dans la monnaie. La profondeur seule ne le justifie jamais, seul le dividende le fait.}
\end{solution}

\begin{solution}{iq:iv2:options-parity-and-greeks:38}
Le put européen vaut environ 6.29~\$, et le put américain environ 7.99~\$. On exerce par anticipation au nœud bas après une période, lorsque l'action vaut 80, et la prime vaut environ 1.70~\$.

La probabilité risque-neutre de hausse vaut $(1.05 - 0.8)/(1.2 - 0.8) = 0.625$. Les gains finaux valent 0 à 144, 4 à 96 et 36 à 64.

Au nœud haut (120), conserver l'option vaut $(0.625 \times 0 + 0.375 \times 4)/1.05 \approx 1.43$, et l'exercer rapporte 0. Au nœud bas (80), la conserver vaut $(0.625 \times 4 + 0.375 \times 36)/1.05 \approx 15.24$, et l'exercer rapporte 20. On exerce donc à ce nœud.

La valeur européenne vaut $(0.625 \times 1.43 + 0.375 \times 15.24)/1.05 \approx 6.29$. La valeur américaine vaut $(0.625 \times 1.43 + 0.375 \times 20)/1.05 \approx 7.99$.

Évaluer une option américaine sur un arbre\index{arbre binomial} revient à faire cette comparaison à chaque nœud.

\iqwrong{Comparer l'exercice avec la valeur finale de l'option. Il faut comparer avec l'espérance actualisée à ce nœud.}
\end{solution}

\begin{solution}{iq:iv2:options-parity-and-greeks:39}
Autour de 100, le gain d'une journée vaut environ 0.092~\$. Autour de 120, il vaut environ 0.011~\$. La même volatilité réalisée donne donc des profits très différents, et le premier dépasse plus de 8 fois le second. La volatilité réalisée seule ne mesure pas votre profit.

Un mouvement de 2~\% à 100 vaut 2~\$, et $\tfrac12 \times 0.046 \times 4 = 0.092$. À 120, il vaut 2.40~\$, et $\tfrac12 \times 0.0038 \times 5.76 \approx 0.011$.

Votre profit le long d'une trajectoire vaut $\sum\tfrac12\Gamma_tS_t^2(r_t^2 - \sigma_i^2\,dt)$, pondéré par le gamma. Les jours où l'action se trouve près du strike comptent donc le plus. C'est donc la volatilité réalisée de ces jours-là qui compte. Un swap de variance\index{swap de variance} est conçu pour supprimer cette dépendance, en détenant un gamma en dollars constant.

\iqwrong{Dire que l'on gagne dès que la volatilité réalisée dépasse la volatilité implicite. On peut observer une volatilité réalisée supérieure à l'implicite et perdre quand même, si les grands mouvements sont survenus alors que le gamma était faible.}
\end{solution}

\begin{solution}{iq:iv2:options-parity-and-greeks:40}
On achète environ 1.135 strangles, on encaisse environ 6.99~\$, et l'on gagne environ 0.10~\$ si la volatilité monte à 25~\%, et environ 0.16~\$ si elle baisse à 15~\%. On est long en volga\index{volga}~: le vega augmente dans le sens du mouvement.

Le vega du straddle vaut environ 79.4 pour 1.00, et celui du strangle environ 69.9, si bien que le rapport vaut 1.135. À 20~\%, le straddle coûte environ 15.93, et 1.135 strangles coûtent environ 8.95. On encaisse donc environ 6.99.

Le vega d'une option à la monnaie varie à peine avec la volatilité, car son prix est presque linéaire en volatilité. Le vega d'une option hors de la monnaie croît avec la volatilité. Cette position a donc un vega nul au départ, et une dérivée seconde positive. Autrement dit, elle profite d'un grand mouvement de la volatilité implicite, dans un sens comme dans l'autre.

\iqwrong{Penser qu'un vega nul signifie l'absence d'exposition à la volatilité. Cela signifie l'absence d'exposition du premier ordre à ce niveau, et c'est précisément le second ordre que la position détient.}
\iqwhy{La volga est ce qu'un teneur de marché vendeur des ailes se fait payer pour porter, et c'est l'une des raisons pour lesquelles les ailes se négocient à une volatilité implicite plus élevée.}
\end{solution}

\begin{solution}{iq:iv2:options-parity-and-greeks:41}
L'écart type vaut environ 0.44~\$, soit environ 5.6~\% de la prime de 7.97~\$. Couvrir quatre fois plus souvent le divise par deux.

L'approximation est $\sqrt{\pi/4}\,\sigma\,\mathcal V/\sqrt N$, où $\mathcal V$ désigne le vega pour 1.00 de volatilité~: $0.886 \times 0.2 \times 39.7/\sqrt{252} \approx 0.443$. Une simulation de 20\,000 trajectoires donne environ 0.44.

L'erreur est proportionnelle à $1/\sqrt N$. Chaque pas laisse un résidu proportionnel au gamma multiplié par le carré du mouvement, diminué de son espérance, et ces erreurs s'ajoutent comme des variances indépendantes.

Un teneur de marché doit donc faire payer ce risque. Il existe même si la volatilité réelle est exactement l'implicite. Couvrir plus souvent le réduit, mais à un coût plus élevé.

\iqwrong{Dire que l'option est «~sans risque une fois couverte~». Cela n'est vrai que pour une couverture continue, et il subsiste toujours un résidu, qui décroît lentement avec la fréquence.}
\end{solution}

\begin{solution}{iq:iv2:options-parity-and-greeks:42}
Porter la volatilité de 20~\% à environ 21.2~\%, ce qui ajoute environ 0.49~\$ à la prime. Une simulation des coûts de rééquilibrage donne environ 0.50~\$.

La règle de Leland\index{règle de Leland} s'écrit $\hat\sigma^2 = \sigma^2(1 + A)$, où $A = \sqrt{2/\pi}\,k/(\sigma\sqrt{\Delta t})$ et $k$ désigne le coût aller-retour. Ici, $\sigma\sqrt{\Delta t} = 0.0126$, et $A = 0.798 \times 0.002/0.0126 \approx 0.127$. D'où $\hat\sigma = 0.2\sqrt{1.127} \approx 21.2\%$. Avec un vega de 39.7 pour 1.00, cela représente environ $0.0123 \times 39.7 \approx 0.49$.

Le nombre espéré d'actions échangées à chaque pas vaut $\Gamma S\sigma\sqrt{\Delta t}\sqrt{2/\pi}$, puisque le mouvement absolu espéré est cette fraction d'un écart type. Le coût de ces transactions, imputé au gamma, ressemble à une volatilité supplémentaire. La valeur simulée omet le premier achat d'environ 54 actions pour 100, qui ajoute environ 0.05~\$.

Un vendeur fait monter la volatilité, et un acheteur la fait baisser, car les coûts jouent contre lui des deux côtés.

\iqwrong{Supposer qu'une couverture plus fréquente aide toujours. Quand $\Delta t$ diminue, $A$ croît, et à un certain point vos coûts l'emportent sur la réduction du risque. Choisir la fréquence est donc déjà le problème.}
\end{solution}

\begin{solution}{iq:iv2:options-parity-and-greeks:43}
Il faut environ 3.45 lots d'un mois, et le gamma net vaut environ $+0.218$ par dollar. Si la volatilité bouge proportionnellement, un seul lot suffit.

Les vegas valent environ 39.7 et 11.5 pour 1.00. D'où $39.7/11.5 \approx 3.45$. Les gammas valent environ 0.0198 et 0.0691. Le net vaut $3.45 \times 0.0691 - 0.0198 \approx 0.218$.

En pratique, la volatilité implicite à court terme bouge davantage. Si un mouvement d'un mois vaut $\sqrt{12}$ fois celui d'un an, un lot de vega d'un mois représente $11.5 \times 3.46 \approx 39.8$ de vega équivalent à un an. Un lot de chaque s'annule donc. C'est l'idée qui sous-tend un vega pondéré\index{vega pondéré}.

\iqwrong{Traiter un calendrier neutre en vega comme neutre tout court. Il ne l'est que face à un mouvement parallèle, et c'est une position très longue en gamma.}
\end{solution}

\begin{solution}{iq:iv2:options-parity-and-greeks:44}
Pas toujours. La variation de valeur par point vaut environ $+0.0070$ à 90, $-0.0020$ à 100 et $-0.0107$ à 110. Le signe change vers 98.02, et non à 100.

La valeur vaut $\Phi(d_2)$, si bien que sa dérivée par rapport à $\sigma$ vaut $\varphi(d_2)\,\partial_\sigma d_2 = -\varphi(d_2)d_1/\sigma$. Elle est nulle lorsque $d_1 = 0$, soit $S = Ke^{-\sigma^2T/2} = 100e^{-0.02} \approx 98.02$.

En dessous de ce niveau, davantage de volatilité donne une meilleure chance d'atteindre le strike. Au-dessus, davantage de volatilité donne une plus grande chance de retomber. Un call ordinaire a un vega positif partout, puisqu'il n'est jamais pénalisé par une volatilité plus forte.

\iqwrong{Supposer que toutes les options profitent de la volatilité. Les gains qui ne sont pas convexes par rapport à l'action peuvent y perdre.}
\end{solution}

\begin{solution}{iq:iv2:options-parity-and-greeks:45}
Si les teneurs de marché sont longs en gamma et que l'action monte de 1~\$, ils vendent 1 million d'actions, soit 20~\% du volume d'une journée, et si elle baisse, ils achètent. S'ils sont courts, ils achètent à la hausse et vendent à la baisse, ce qui amplifie le mouvement. Un gamma long tend donc à amortir les mouvements et à figer le cours, tandis qu'un gamma court tend à les amplifier.

L'estimation n'est pas fiable. Elle suppose que les teneurs de marché se trouvent d'un seul côté de toutes les transactions des clients, et elle ignore la manière dont ils se sont couverts. Les positions ouvertes ne disent pas non plus qui a acheté et qui a vendu. Il faut donc confronter l'affirmation aux données de volatilité réalisée effective, et la traiter comme un a priori faible.

\iqwrong{Présenter cela comme un fait établi du marché. C'est un mécanisme appuyé par des observations, mais qui repose sur une estimation que l'on ne voit pas directement.}
\end{solution}

\begin{solution}{iq:iv2:options-parity-and-greeks:46}
Dériver $C(\lambda S, \lambda K) = \lambda C(S, K)$ par rapport à $\lambda$ en 1. On obtient $C = S\partial_S C + K\partial_K C$. Le delta vaut donc $(C - K\partial_K C)/S$. D'après les cotations, $\partial_K C \approx (7.52 - 8.44)/2 = -0.46$. Le delta vaut donc $(7.97 + 100 \times 0.46)/100 = 0.5397$, contre $\Phi(0.1) \approx 0.5398$ dans Black-Scholes.

Cela suppose que le prix du call ne dépend que du rapport entre l'action et le strike, si bien que le smile se déplace avec l'action~: c'est le sticky moneyness\index{sticky moneyness}. Le delta se lit alors dans la pente du call par rapport au strike, sans recourir à un modèle. Si en revanche le smile reste fixe par strike (sticky strike), le delta est un autre nombre, le skew venant s'y ajouter ou en être retranché.

\iqwrong{Penser qu'il existe un seul delta correct. Le modèle qui produit le delta contient une hypothèse sur la manière dont le smile se déplace.}
\end{solution}

\begin{solution}{iq:iv2:options-parity-and-greeks:47}
Tous les prix décroissent avec le strike, et aucun spread ne dépasse sa largeur. La condition de convexité est violée en 100, 110 et 120. Achetez donc le 100 et le 120, et vendez deux fois le 110. On vous paie 1~\$, et votre pire issue est ce gain de 1~\$, tandis que la meilleure vaut 11~\$. La table implique des probabilités d'environ 60~\%, 25~\% et 35~\%, et la hausse de 25~\% à 35~\% est impossible.

Le butterfly coûte $8 - 2 \times 5.5 + 2 = -1$. Son gain vaut 0 en dessous de 100 ou égal à 100, et au-delà de 120, et il atteint 10 en 110. D'où $-(-1) + 0 = 1$ dans le pire cas, et $1 + 10 = 11$ dans le meilleur.

Pentes~: $(14 - 8)/10 = 0.60$, $(8 - 5.5)/10 = 0.25$, $(5.5 - 2)/10 = 0.35$. La valeur d'un call spread par unité de largeur vaut à peu près la probabilité de finir au-dessus de son milieu, et celle-ci doit décroître avec le strike. C'est donc la même infraction que celle du butterfly.

\iqwrong{Vérifier seulement que les prix décroissent avec le strike. La convexité est la condition la plus souvent violée.}
\end{solution}

\begin{solution}{iq:iv2:options-parity-and-greeks:48}
Toute fonction $f$ deux fois dérivable vérifie
\[
f(S) = f(K_0) + f'(K_0)(S - K_0) + \int_0^{K_0}f^{\prime\prime}(K)(K - S)^+\,dK + \int_{K_0}^\infty f^{\prime\prime}(K)(S - K)^+\,dK.
\]
Pour $S_T^2$, $f^{\prime\prime} = 2$ pour chaque strike. On détient donc 2 puts de chaque strike inférieur à 100 et 2 calls de chaque strike supérieur à 100. Avec des taux nuls, le strip vaut environ 408.11. Avec l'obligation de 10\,000 et la valeur nulle du forward, le total s'élève à 10\,408.11, soit exactement $S^2e^{\sigma^2T}$.

La formule découle du théorème de Taylor avec reste intégral. Une intégration par parties transforme les termes de dérivée seconde en options. Il s'agit de la réplication statique de Carr-Madan\index{Carr-Madan replication}.

Pour $f = S^2$~: $f(K_0) = 10\,000$, et $f'(K_0)(S - K_0)$ est un forward de valeur nulle. Les options sont pondérées par $f^{\prime\prime}(K)\,dK = 2\,dK$.

Pour vérifier, Black-Scholes donne $E[S_T^2] = S_0^2e^{\sigma^2T} = 10\,000e^{0.04} \approx 10\,408.11$. L'intégration numérique du strip confirme ce résultat.

Tout payoff européen possède donc une valeur indépendante du modèle, à partir des prix des options, sans qu'il faille modéliser la trajectoire. C'est le fondement du prix du swap de variance, obtenu en prenant $f = \ln$.

\iqwrong{Croire qu'il faut un modèle de volatilité pour valoriser un payoff régulier. Le smile le valorise, et la seule question est de savoir combien de strikes on peut négocier.}
\end{solution}

\begin{solution}{iq:iv2:options-parity-and-greeks:49}
Votre profit espéré vaut environ 1.99~\$, soit la différence entre l'option à 30~\% et celle à 20~\%. Une couverture par les deltas implicites donne un profit le long d'une trajectoire de $\int\tfrac12\Gamma_iS^2(\sigma_i^2 - \sigma_r^2)\,dt$, avec une moyenne de 1.99 pour une couverture quotidienne et un écart type d'environ 0.74. Une couverture par les deltas réalisés donne un profit certain de $V(\sigma_i) - V(\sigma_r)$ en couverture continue, et d'environ 1.99 avec un écart type de 0.44 en couverture quotidienne. Si l'on connaissait les 20~\%, on couvrirait avec ce chiffre. La plupart des desks utilisent la volatilité implicite, puisque personne ne connaît la réalisée, et le P\&L s'accumule alors régulièrement à un taux connu.

Couvrez à $\sigma_r$. Soit $V_r(S, t)$ la valeur de l'option à la volatilité $\sigma_r$. D'après Itô, sa variation vaut $\Delta_r\,dS$, le thêta et le gamma se compensant à cette volatilité réalisée. La couverture compense donc l'option dès le départ, à une valeur de $V_r$, et vous conservez $V_i - V_r$ sur chaque trajectoire. Le 0.44 ne tient qu'à la discrétisation.

Couvrez à $\sigma_i$. Appliquez Itô à $V_i$ sous la dynamique réalisée~: $dV_i - \Delta_i\,dS = (\theta_i + \tfrac12\sigma_r^2S^2\Gamma_i)dt = \tfrac12\Gamma_iS^2(\sigma_r^2 - \sigma_i^2)dt$. Pour une option vendue, vous conservez l'opposé de ce terme. En prenant l'espérance, on retrouve le même $V_i - V_r$.

\iqwrong{Affirmer que la couverture à l'implicite est «~la bonne~» parce que c'est la volatilité du marché. Elle ne réduit pas la variance, puisque la trajectoire réalisée détermine votre P\&L. Son avantage réside dans le calendrier du P\&L.}
\end{solution}

\begin{solution}{iq:iv2:options-parity-and-greeks:50}
Le niveau d'exercice vaut environ 71.43~\$, et le put vaut environ 12.32~\$ lorsque l'action est à 100. Quand le taux tend vers 0, on n'exerce jamais, et la valeur tend vers le strike.

Une option perpétuelle ne dépend pas du temps, si bien que l'équation de Black-Scholes devient $\tfrac12\sigma^2S^2V^{\prime\prime} + rSV' - rV = 0$. Essayons $V = S^p$~: $\tfrac12\sigma^2p(p - 1) + rp - r = 0$, avec pour racines $p = 1$ et $p = -\gamma$, où $\gamma = 2r/\sigma^2 = 2.5$. La valeur doit s'annuler pour $S$ grand, donc $V = AS^{-\gamma}$.

Au niveau d'exercice $S^*$, on a $V(S^*) = K - S^*$. Pour un niveau $L$ quelconque, la valeur en $S$ vaut alors $(K - L)(S/L)^{-\gamma}$. Pour maximiser sur $L$, on annule la dérivée de $(K - L)L^\gamma$, ce qui donne $L = \gamma K/(1 + \gamma) = 71.43$. Ce choix coïncide avec le recollement C1~: $V'(S^*) = -1$.

Avec l'action à 100~: $28.57 \times (1.4)^{-2.5} \approx 12.32$.

Quand $r \to 0$, $\gamma \to 0$, donc $L \to 0$, et $V \to K$. Sans intérêt à gagner sur la trésorerie, on attend. À dérive nulle, l'action tend vers zéro avec probabilité un, et la valeur du put tend vers la totalité du strike.

\iqwrong{Fixer le niveau d'exercice en égalant la valeur à celle d'une option européenne. Le niveau découle de la maximisation sur le choix de l'exercice, et c'est l'essence de la valorisation américaine.}
\iqfollow{Que fait une échéance finie~?}{Il n'existe alors aucune forme fermée. Un arbre ou une grille aux différences finies la résout, et la frontière remonte vers le strike à mesure que le temps passe.}
\end{solution}
